\subsubsection{伪单调算子}\label{subsec:pseudo-monotone-operators}

本节考察一类稍为一般的非线性映射, 称为\emph{伪单调算子}. 在应用中, 前一节施加的假设往往强得没有必要. 特别地, 单调性假设 \ref{hyp:monotone-pde-monotonicity} 同时涉及一阶导数和函数本身. 正如本章将看到的, 实际上只需对最高阶导数作单调性假设: 紧性将处理低阶项.

\begin{Definition}
\label{def:pseudo-monotone-operator}
设 $X$ 为自反 Banach 空间. 若 $T$ 有界, 且每当
\begin{equation}\label{eq:pseudo-monotonicity-weak-sequence}
u_j\rightharpoonup\overline u\quad\text{in }X
\end{equation}
和
\begin{equation}\label{eq:pseudo-monotonicity-limsup-hypothesis}
\limsup_{j\to\infty}(T(u_j),u_j-\overline u)\leq0,
\end{equation}
成立时, 都有
\begin{equation}\label{eq:pseudo-monotonicity-liminf-conclusion}
\liminf_{j\to\infty}(T(u_j),u_j-v)\geq(T(\overline u),\overline u-v)\quad\text{for all }v\in X.
\end{equation}
则称算子 $T:X\to X^*$ 为\textbf{伪单调的}.
\end{Definition}

只需对 Browder-Minty 定理的证明略作修改, 就可证明下面的结果.

\begin{Theorem}
\label{thm:pseudo-monotone-surjectivity}
设 $X$ 为实自反 Banach 空间, 并假设 $T:X\to X^*$ 连续, 强制且伪单调. 则对每个 $g\in X^*$, 方程
\begin{equation}\label{eq:pseudo-monotone-surjectivity-equation}
T(u)=g.
\end{equation}
存在解 $u\in X$.
\end{Theorem}

证明留给读者 (Problem~\ref{ex:pseudo-monotone-surjectivity-proof}).

实际中, 下面的条件比伪单调性更容易验证.

\begin{Definition}
\label{def:calculus-variations-type-operator}
设 $X$ 为自反 Banach 空间. 若算子 $T:X\to X^*$ 有界, 并具有表示
\begin{equation}\label{eq:calculus-variations-type-diagonal-representation}
T(u)=\widehat T(u,u)
\end{equation}
其中映射
\begin{equation}\label{eq:calculus-variations-type-two-argument-mapping}
X\times X\ni(u,v)\mapsto\widehat T(u,v)\in X^*
\end{equation}
满足下列假设, 则称它为\textbf{变分法型}的.\begin{enumerate}[label=CV-\arabic*.,ref=CV-\arabic*]\item\label{hyp:cv-diagonal-monotonicity} 对每个 $u\in X$, 映射 $v\mapsto\widehat T(u,v)$ 从 $X$ 到 $X^*$ 有界且连续, 并且
\begin{equation}\label{eq:calculus-variations-type-diagonal-monotonicity}
(\widehat T(u,u)-\widehat T(u,v),u-v)\geq0\quad\text{for all }v\in X.
\end{equation}
\item\label{hyp:cv-first-variable-continuity} 对每个 $v\in X$, 映射 $u\mapsto\widehat T(u,v)$ 从 $X$ 到 $X^*$ 有界且连续.\item\label{hyp:cv-vanishing-defect-weak-continuity} 若
\begin{equation}\label{eq:calculus-variations-type-cv-three-weak-sequence}
u_j\rightharpoonup\overline u\text{ in }X
\end{equation}
且
\begin{equation}\label{eq:calculus-variations-type-vanishing-diagonal-pairing}
(\widehat T(u_j,u_j)-\widehat T(u_j,\overline u),u_j-\overline u)\to0,
\end{equation}
则对每个 $v\in X$
\begin{equation}\label{eq:calculus-variations-type-weak-image-continuity}
\widehat T(u_j,v)\rightharpoonup\widehat T(\overline u,v)\text{ in }X^*.
\end{equation}
\item\label{hyp:cv-weak-duality-pairing} 若
\begin{equation}\label{eq:calculus-variations-type-cv-four-weak-sequence}
u_j\rightharpoonup\overline u\text{ in }X
\end{equation}
且
\begin{equation}\label{eq:calculus-variations-type-cv-four-weak-image}
\widehat T(u_j,v)\rightharpoonup\psi\text{ in }X^*,
\end{equation}
则
\begin{equation}\label{eq:calculus-variations-type-cv-four-pairing-convergence}
(\widehat T(u_j,v),u_j)\to(\psi,\overline u).
\end{equation}
\end{enumerate}
\end{Definition}

如上面所指出的, 有下面的结果.

\begin{Theorem}
\label{thm:cv-type-implies-pseudo-monotonicity}
若 $T$ 为变分法型的, 则 $T$ 伪单调.
\end{Theorem}

\begingroup
\renewcommand{\endProof}{\par\nobreak\hfill$\blacksquare$\par\vspace{0.1cm}}
\begin{Proof}
设 $u_j\rightharpoonup\overline u$ 在 $X$ 中成立, 并假设
\begin{equation}\label{eq:cv-pseudo-proof-limsup-assumption}
\limsup_{j\to\infty}(T(u_j),u_j-\overline u)\leq0.
\end{equation}
要证明
\begin{equation}\label{eq:cv-pseudo-proof-liminf-goal}
\liminf_{j\to\infty}(T(u_j),u_j-v)\geq(T(\overline u),\overline u-v)\quad\text{for every }v\in X.
\end{equation}
由于 $\widehat T(u_j,\overline u)$ 在 $X^*$ 中有界, 可以抽取子序列 $u_j$, 使对某个 $\psi\in X^*$ 有
\begin{equation}\label{eq:cv-pseudo-proof-bounded-weak-image-subsequence}
\widehat T(u_j,\overline u)\rightharpoonup\psi\quad\text{in }X^*,
\end{equation}
现在利用 \ref{hyp:cv-weak-duality-pairing}, 得到
\begin{equation}\label{eq:cv-pseudo-proof-duality-pair-limit}
\lim_{j\to\infty}(\widehat T(u_j,\overline u),u_j)=(\psi,\overline u).
\end{equation}
因此, 若定义
\begin{equation}\label{eq:cv-pseudo-proof-nonnegative-defect-definition}
x_j:=(\widehat T(u_j,u_j)-\widehat T(u_j,\overline u),u_j-\overline u)\in\mathbb{R},
\end{equation}
则有
\begin{equation}\label{eq:cv-pseudo-proof-defect-limsup-bound}
\begin{aligned}\limsup_{j\to\infty}x_j=\limsup_{j\to\infty}\bigl[&(T(u_j),u_j-\overline u)-(\widehat T(u_j,\overline u),u_j)\\&+(\widehat T(u_j,\overline u),\overline u)\bigr]\leq0.\end{aligned}
\end{equation}
这里使用了 \eqref{eq:cv-pseudo-proof-limsup-assumption}, \eqref{eq:cv-pseudo-proof-bounded-weak-image-subsequence} 和 \eqref{eq:cv-pseudo-proof-duality-pair-limit}. 由于 \ref{hyp:cv-diagonal-monotonicity} 蕴含 $x_j\geq0$, 得到
\begin{equation}\label{eq:cv-pseudo-proof-defect-vanishing}
\lim_{j\to\infty}x_j=0.
\end{equation}
因此可以利用 \ref{hyp:cv-vanishing-defect-weak-continuity}, 得到
\begin{equation}\label{eq:cv-pseudo-proof-test-image-limit}
\widehat T(u_j,v)\rightharpoonup\widehat T(\overline u,v)\text{ in }X^*\quad\text{for all }v\in X.
\end{equation}
从而再次利用 \ref{hyp:cv-weak-duality-pairing}, 得到\[(\widehat T(u_j,v),u_j)\to(\widehat T(\overline u,v),\overline u)\]即
\begin{equation}\label{eq:cv-pseudo-proof-vanishing-test-pairing}
(\widehat T(u_j,v),u_j-\overline u)\to0\quad\text{for all }v\in X.
\end{equation}
现在利用这一事实和 $x_j\geq0$, 得到
\begin{equation}\label{eq:cv-pseudo-proof-diagonal-comparison}
(T(u_j),u_j-\overline u)\geq(\widehat T(u_j,\overline u),u_j-\overline u)\to0.
\end{equation}
结合 \eqref{eq:cv-pseudo-proof-limsup-assumption}, 得到
\begin{equation}\label{eq:cv-pseudo-proof-diagonal-pairing-limit}
(T(u_j),u_j-\overline u)\to0.
\end{equation}
现在取 \ref{hyp:cv-diagonal-monotonicity} 中的不等式\[(T(u_j)-\widehat T(u_j,w),u_j-w)\geq0\quad\text{for all }w\in X\]并代入
\begin{equation}\label{eq:cv-pseudo-proof-intermediate-test-point}
w=(1-\theta)\overline u+\theta v,
\end{equation}
其中 $\theta\in(0,1)$. 得到
\begin{equation}\label{eq:cv-pseudo-proof-theta-inequality}
\begin{aligned}\theta(T(u_j),\overline u-v)\geq&-(T(u_j),u_j-\overline u)+(\widehat T(u_j,w),u_j-\overline u)\\&+\theta(\widehat T(u_j,w),\overline u-v).\end{aligned}
\end{equation}
将其除以 $\theta$, 并利用 \eqref{eq:cv-pseudo-proof-vanishing-test-pairing} 和 \eqref{eq:cv-pseudo-proof-diagonal-pairing-limit}, 得到\[\begin{aligned}
\liminf_{j\to\infty}(T(u_j),u_j-v)&=\liminf_{j\to\infty}(T(u_j),u_j-\overline u)+\liminf_{j\to\infty}(T(u_j),\overline u-v)\\
&\geq\liminf_{j\to\infty}(\widehat T(u_j,w),\overline u-v)\\
&=(\widehat T(\overline u,w),\overline u-v)\\
&=(\widehat T(\overline u,(1-\theta)\overline u+\theta v),\overline u-v).
\end{aligned}\]令 $\theta\searrow0$, 得到
\begin{equation}\label{eq:cv-pseudo-proof-liminf-conclusion}
\liminf_{j\to\infty}(T(u_j),u_j-v)\geq(T(\overline u),\overline u-v)\quad\text{for all }v\in X.
\end{equation}
由于上述论证对原序列的任意子序列都成立, 不等式 \eqref{eq:cv-pseudo-proof-liminf-conclusion} 对整个原序列成立. 证明完成.
\end{Proof}
\endgroup

下面的结论立即由前面的结果得到.

\begin{Corollary}
\label{cor:cv-type-surjectivity}
设 $X$ 为实自反 Banach 空间, 并假设 $T:X\to X^*$ 连续, 强制且为变分法型的. 则对每个 $g\in X^*$, 方程
\begin{equation}\label{eq:calculus-variations-type-surjectivity-equation}
T(u)=g.
\end{equation}
存在解 $u\in X$.
\end{Corollary}

\subsubsection{对 PDE 的应用}\label{subsec:cv-type-application-to-pdes}

设 $\Omega\subset\mathbb{R}^n$ 为具有光滑边界的有界区域. 考虑如下形式的拟线性二阶微分算子
\begin{equation}\label{eq:cv-quasilinear-second-order-operator}
\widetilde A(u)(\bm{x})=-\sum_{i=1}^n\dfrac{\partial}{\partial x_i}a_i(\bm{x},u(\bm{x}),\nabla u(\bm{x}))+a_0(\bm{x},u(\bm{x}),\nabla u(\bm{x})).
\end{equation}
目标是对适当的 $f$ 求解方程
\begin{equation}\label{eq:cv-quasilinear-dirichlet-equation}
\widetilde A(u)=f
\end{equation}
的 Dirichlet 问题. 形式地, 定义二元型
\begin{equation}\label{eq:cv-quasilinear-bivariate-form}
\begin{aligned}\widetilde B(u,v):=\int_\Omega&\sum_{i=1}^n a_i(\bm{x},u(\bm{x}),\nabla u(\bm{x}))\dfrac{\partial v(\bm{x})}{\partial x_i}\\&+a_0(\bm{x},u(\bm{x}),\nabla u(\bm{x}))v(\bm{x})\,\mathrm{d}\bm{x}.\end{aligned}
\end{equation}
对函数
\begin{equation}\label{eq:cv-pde-coefficient-functions-domain}
\Omega\times\mathbb{R}\times\mathbb{R}^n\ni(\bm{x},\eta,\bm{x}ii)\mapsto a_i(\bm{x},\eta,\bm{x}ii)\in\mathbb{R},\quad i=0,\ldots,n.
\end{equation}
作如下假设.

\begin{enumerate}[label=HP-\arabic*.,ref=HP-\arabic*]\item\label{hyp:cv-pde-spatial-continuity} 对每个 $i=0,\ldots,n$,
\begin{equation}\label{eq:cv-pde-coefficient-spatial-continuity}
\bm{x}\mapsto a_i(\bm{x},\eta,\bm{x}ii)
\end{equation}
对每个固定的 $(\eta,\bm{x}ii)\in\mathbb{R}^{n+1}$ 都属于 $C_b(\overline\Omega)$.

\item\label{hyp:cv-pde-value-continuity} 对每个 $i=0,\ldots,n$,
\begin{equation}\label{eq:cv-pde-coefficient-value-gradient-continuity}
(\eta,\bm{x}ii)\mapsto a_i(\bm{x},\eta,\bm{x}ii)
\end{equation}
对每个 $\bm{x}\in\Omega$ 都属于 $C(\mathbb{R}^{n+1})$.

\item\label{hyp:cv-pde-growth} 存在 $p\in(1,\infty)$, 常数 $c_0>0$ 和函数 $k\in L^q(\Omega)$ ($\frac{1}{p}+\frac{1}{q}=1$), 使对每个 $\bm{x}\in\Omega$ 以及每个 $(\eta,\bm{x}ii)\in\mathbb{R}^{n+1}$, 都有
\begin{equation}\label{eq:cv-pde-coefficient-growth-bound}
|a_i(\bm{x},\eta,\bm{x}ii)|\leq c[|\eta|^{p-1}+|\bm{x}ii|^{p-1}+k(\bm{x})],
\end{equation}
其中 $i=0,\ldots,n$ 任意.

\item\label{hyp:cv-pde-strict-gradient-monotonicity} 对每个满足 $\bm{x}ii\neq\bm{x}ii^*$ 的 $\bm{x}ii\in\mathbb{R}^n$ 和 $\bm{x}ii^*\in\mathbb{R}^n$, 以及每个 $\eta\in\mathbb{R}$ 和 $\bm{x}\in\Omega$, 都有
\begin{equation}\label{eq:cv-pde-strict-highest-order-monotonicity}
\sum_{i=1}^n[a_i(\bm{x},\eta,\bm{x}ii)-a_i(\bm{x},\eta,\bm{x}ii^*)](\bm{x}i_i-\bm{x}i_i^*)>0.
\end{equation}


\item\label{hyp:cv-pde-form-coercivity}
\begin{equation}\label{eq:cv-pde-absolute-bivariate-coercivity-hypothesis}
\dfrac{|\widetilde B(v,v)|}{\|v\|_{1,p}}\to\infty\quad\text{as }\|v\|_{1,p}\to\infty.
\end{equation}


\item\label{hyp:cv-pde-uniform-gradient-coercivity} 对每个 $\bm{x}\in\Omega$, 且对有界集中的 $|\eta|$ 一致地, 有
\begin{equation}\label{eq:cv-pde-uniform-gradient-coercivity-hypothesis}
\sum_{i=1}^na_i(\bm{x},\eta,\bm{x}ii)\bm{x}i_i\dfrac{1}{|\bm{x}ii|+|\bm{x}ii|^{p-1}}\to\infty\quad\text{as }|\bm{x}ii|\to\infty.
\end{equation}
\end{enumerate}

注意, 由假设 \ref{hyp:cv-pde-spatial-continuity}, \ref{hyp:cv-pde-value-continuity} 和 \ref{hyp:cv-pde-growth}, 二元型 $\widetilde B(u,v)$ 对属于 $W^{1,p}(\Omega)$ 的 $u,v$ 有良好定义, 其中 $p\in(1,\infty)$ 如 \ref{hyp:cv-pde-growth} 中所给定. 设 $f\in W^{-1,q}(\Omega)$. 与前面相同, 若
\begin{equation}\label{eq:cv-quasilinear-dirichlet-weak-solution}
\widetilde B(u,v)=(f,v)\quad\text{for every }v\in W_0^{1,p}(\Omega).
\end{equation}
则称 $u\in W_0^{1,p}(\Omega)$ 为 \eqref{eq:cv-quasilinear-dirichlet-equation} 的 Dirichlet 问题的弱解. 假设 \ref{hyp:cv-pde-spatial-continuity}, \ref{hyp:cv-pde-value-continuity} 和 \ref{hyp:cv-pde-growth} 还蕴含, 对每个固定的 $u\in W_0^{1,p}(\Omega)$, 映射
\begin{equation}\label{eq:cv-quasilinear-bounded-test-functional}
W_0^{1,p}(\Omega)\ni w\mapsto\widetilde B(u,w)\in\mathbb{R}
\end{equation}
为有界线性泛函. 因此, 存在映射
\begin{equation}\label{eq:cv-quasilinear-induced-dual-mapping}
W_0^{1,p}(\Omega)\ni u\mapsto\overline T(u)\in W^{-1,q}(\Omega)
\end{equation}
使得
\begin{equation}\label{eq:cv-quasilinear-induced-pairing}
\widetilde B(u,w)=(\overline T(u),w)\quad\text{for all }w\in W_0^{1,p}(\Omega).
\end{equation}


现在证明下面的结论.

\begin{Theorem}
\label{thm:cv-pde-operator-is-variations-type}
设 $a_i$ 满足假设 \ref{hyp:cv-pde-spatial-continuity}--\ref{hyp:cv-pde-uniform-gradient-coercivity}. 则算子 $\overline T$ 为变分法型的.
\end{Theorem}

\begingroup
\renewcommand{\endProof}{\par}
\begin{Proof}
当然, 这里的目标之一是将高阶导数与低阶导数的作用分开; 因此定义
\begin{equation}\label{eq:cv-quasilinear-three-argument-form-splitting}
\widetilde B(u,v,w):=\mathcal{B}_1(u,v,w)+\mathcal{B}_0(u,w),
\end{equation}
其中
\begin{equation}\label{eq:cv-quasilinear-highest-order-three-argument-form}
\mathcal{B}_1(u,v,w):=\sum_{i=1}^n\int_\Omega a_i(\bm{x},u(\bm{x}),\nabla v(\bm{x}))\dfrac{\partial w}{\partial x_i}(\bm{x})\,\mathrm{d}\bm{x},
\end{equation}

\begin{equation}\label{eq:cv-quasilinear-lower-order-form}
\mathcal{B}_0(u,w):=\int_\Omega a_0(\bm{x},u(\bm{x}),\nabla u(\bm{x}))w(\bm{x})\,\mathrm{d}\bm{x}.
\end{equation}
利用与前面相同的论证, 可知存在映射
\begin{equation}\label{eq:cv-quasilinear-two-argument-dual-mapping}
W_0^{1,p}(\Omega)\times W_0^{1,p}(\Omega)\ni(u,v)\mapsto\widehat T(u,v)\in W^{-1,q}(\Omega)
\end{equation}
使得
\begin{equation}\label{eq:cv-quasilinear-three-argument-duality-pairing}
\widetilde B(u,v,w)=(\widehat T(u,v),w)\quad\text{for all }w\in W_0^{1,p}(\Omega).
\end{equation}
此外, 有
\begin{equation}\label{eq:cv-quasilinear-diagonal-operator-representation}
\overline T(u)=\widehat T(u,u).
\end{equation}
现在注意到, 关于 Nemytskii 算子的结果以及假设 \ref{hyp:cv-pde-spatial-continuity}, \ref{hyp:cv-pde-value-continuity} 和 \ref{hyp:cv-pde-growth} 立即蕴含下面的引理.
\end{Proof}
\endgroup

\begin{Lemma}
\label{lem:cv-pde-nemytskii-partial-argument-continuity}
对每个 $i=0,\ldots,n$, 映射
\begin{equation}\label{eq:cv-pde-nemytskii-full-argument-continuity}
W^{1,p}(\Omega)\ni u\mapsto a_i(\bm{x},u(\bm{x}),\nabla u(\bm{x}))\in L^q(\Omega)
\end{equation}
有界且连续. 此外, 对固定的 $v\in L^p(\Omega)$, 映射
\begin{equation}\label{eq:cv-pde-nemytskii-fixed-value-continuity}
W^{1,p}(\Omega)\ni u\mapsto a_i(\bm{x},v(\bm{x}),\nabla u(\bm{x}))\in L^q(\Omega)
\end{equation}
有界且连续; 对固定的 $w\in W^{1,p}(\Omega)$, 映射
\begin{equation}\label{eq:cv-pde-nemytskii-fixed-gradient-continuity}
L^p(\Omega)\ni v\mapsto a_i(\bm{x},v(\bm{x}),\nabla w(\bm{x}))\in L^q(\Omega)
\end{equation}
有界且连续.
\end{Lemma}

由此得到下面的推论.

\begin{Corollary}
\label{cor:cv-pde-partial-argument-continuity}
有下列结论.\begin{enumerate}\item 算子 $\overline T$ 有界且连续.\item 对每个 $u\in W_0^{1,p}(\Omega)$, 映射
\begin{equation}\label{eq:cv-pde-two-argument-second-variable-continuity}
W_0^{1,p}(\Omega)\ni v\mapsto\widehat T(u,v)\in W^{-1,q}(\Omega)
\end{equation}
有界且连续.\item 对每个 $v\in W_0^{1,p}(\Omega)$, 映射
\begin{equation}\label{eq:cv-pde-two-argument-first-variable-continuity}
W_0^{1,p}(\Omega)\ni u\mapsto\widehat T(u,v)\in W^{-1,q}(\Omega)
\end{equation}
有界且连续.\end{enumerate}
\end{Corollary}

\begingroup
\renewcommand{\endProof}{\par\nobreak\hfill$\blacksquare$\par\vspace{0.1cm}}
\begingroup
\renewenvironment{Proof}{\par}{\par\nobreak\hfill$\blacksquare$\par\vspace{0.1cm}}
\begin{Proof}
这一结果和假设 \ref{hyp:cv-pde-strict-gradient-monotonicity} 表明, 条件 \ref{hyp:cv-diagonal-monotonicity} 和 \ref{hyp:cv-first-variable-continuity} 满足. 因此, 为证明 $\overline T$ 为变分法型的并完成 Theorem~\ref{thm:cv-pde-operator-is-variations-type} 的证明, 只需验证条件 \ref{hyp:cv-vanishing-defect-weak-continuity} 和 \ref{hyp:cv-weak-duality-pairing} 满足. 为检查条件 \ref{hyp:cv-vanishing-defect-weak-continuity}, 假设 $u_j\rightharpoonup\overline u$ 在 $W_0^{1,p}(\Omega)$ 中成立, 且
\begin{equation}\label{eq:cv-pde-vanishing-gradient-monotonicity-defect}
\begin{aligned}
&(\widehat T(u_j,u_j)-\widehat T(u_j,\overline u),u_j-\overline u)\\
&\quad=\mathcal{B}_1(u_j,u_j,u_j-\overline u)-\mathcal{B}_1(u_j,\overline u,u_j-\overline u)\\
&\quad=\int_\Omega\sum_{i=1}^n[a_i(\bm{x},u_j,\nabla u_j)-a_i(\bm{x},u_j,\nabla\overline u)]\dfrac{\partial}{\partial x_i}(u_j-\overline u)\,\mathrm{d}\bm{x}\\
&\quad\to0,
\end{aligned}
\end{equation}
需要证明
\begin{equation}\label{eq:cv-pde-source-unhatted-test-operator-limit}
\begin{aligned}(T(u_j,v),w)&=\mathcal{B}_1(u_j,v,w)+\mathcal{B}_0(u_j,w)\\&\to\mathcal{B}_1(\overline u,v,w)+\mathcal{B}_0(\overline u,w)\\&=(T(\overline u,v),w)\end{aligned}
\end{equation}
对所有 $w\in W_0^{1,p}(\Omega)$ 成立. 现在由紧嵌入, $u_j\to u$ 在 $L^p(\Omega)$ 中强收敛. 因此, 即使不使用假设 \eqref{eq:cv-pde-vanishing-gradient-monotonicity-defect}, 也有\[\begin{aligned}
\mathcal{B}_1(u_j,v,w)&=\int_\Omega\sum_{i=1}^na_i(\bm{x},u_j(\bm{x}),\nabla v(\bm{x}))\dfrac{\partial w}{\partial x_i}(\bm{x})\,\mathrm{d}\bm{x}\\
&\to\mathcal{B}_1(\overline u,v,w).
\end{aligned}\]这里使用了关于 Nemytskii 算子的结果. 因而只需证明
\begin{equation}\label{eq:cv-pde-lower-order-limit-to-verify}
\int_\Omega a_0(\bm{x},u_j(\bm{x}),\nabla u_j(\bm{x}))w(\bm{x})\,\mathrm{d}\bm{x}\to\int_\Omega a_0(\bm{x},\overline u(\bm{x}),\nabla\overline u(\bm{x}))w(\bm{x})\,\mathrm{d}\bm{x}.
\end{equation}
一般情形下的证明是测度论中的一个练习, 这里略去. 请读者参阅 \cite[Lemma 2.2, p.~184]{Li}. 在下面的例子中, 使用如下形式的低阶项
\begin{equation}\label{eq:cv-pde-linear-drift-p-power-lower-order-term}
a_0(\bm{x},u(\bm{x}),\nabla u(\bm{x}))=\bm{b}_1(\bm{x})\cdot\nabla u(\bm{x})+\alpha_0(\bm{x})|u(\bm{x})|^{p-2}u(\bm{x}),
\end{equation}
其中 $\bm{b}_1$ 和 $\alpha_0$ 为有界连续函数. 对这样的项, 条件 \eqref{eq:cv-pde-lower-order-limit-to-verify} 可直接验证 (即不必使用 \eqref{eq:cv-pde-vanishing-gradient-monotonicity-defect}, 不过这一条件对一般证明是必需的).

为检查条件 \ref{hyp:cv-weak-duality-pairing}, 假设 $u_j\rightharpoonup\overline u$ 在 $W_0^{1,p}(\Omega)$ 中成立, 且 $\widehat T(u_j,v)\rightharpoonup\psi$ 在 $W^{-1,q}(\Omega)$ 中成立; 即
\begin{equation}\label{eq:cv-pde-weak-test-image-limit-assumption}
(\widehat T(u_j,v),w)=\mathcal{B}_1(u_j,v,w)+\mathcal{B}_0(u_j,w)\to(\psi,w)
\end{equation}
对每个 $w\in W_0^{1,p}(\Omega)$ 成立. 要证明
\begin{equation}\label{eq:cv-pde-weak-duality-pairing-limit-goal}
(\widehat T(u_j,v),u_j)=\mathcal{B}_1(u_j,v,u_j)+\mathcal{B}_0(u_j,u_j)\to(\psi,\overline u).
\end{equation}
为此, 写成
\begin{equation}\label{eq:cv-pde-weak-duality-pairing-splitting}
(\widehat T(u_j,v),u_j)=(\widehat T(u_j,v),\overline u)+(\widehat T(u_j,v),u_j-\overline u).
\end{equation}
因此, 由 \eqref{eq:cv-pde-weak-test-image-limit-assumption}, 只需证明
\begin{equation}\label{eq:cv-pde-vanishing-weak-test-pairing-exercise}
\lim_{j\to\infty}(\widehat T(u_j,v),u_j-\overline u)=0.
\end{equation}
证明留给读者 (Problem~\ref{ex:cv-pde-vanishing-weak-test-pairing}). 这表明所考虑的算子为变分法型的, 从而完成证明.
\end{Proof}
\endgroup
\endgroup

为得到关于偏微分方程的主要结果, 还需证明一个引理.

\begin{Lemma}
\label{lem:cv-pde-induced-operator-coercivity}
算子 $\overline T$ 为强制的.
\end{Lemma}

证明留给读者 (Problem~\ref{ex:cv-pde-induced-operator-coercivity-proof}). 前述结果最终给出下面关于拟线性椭圆偏微分方程的存在性定理.

\begin{Theorem}
\label{thm:cv-quasilinear-dirichlet-existence}
设 $a_i$ 满足假设 (\ref{hyp:cv-pde-spatial-continuity}--\ref{hyp:cv-pde-uniform-gradient-coercivity}), 并设 $p$ 和 $q$ 如假设 \ref{hyp:cv-pde-growth} 中所定义. 则对每个 $f\in W^{-1,q}(\Omega)$, 拟线性 PDE \eqref{eq:cv-quasilinear-dirichlet-equation} 的 Dirichlet 问题存在弱解 $u\in W_0^{1,p}(\Omega)$.
\end{Theorem}

\begin{Example}
\label{ex:cv-quasilinear-anisotropic-drift-operator}
考虑二阶非线性偏微分算子
\begin{equation}\label{eq:cv-quasilinear-anisotropic-drift-example}
\begin{aligned}\widetilde A(u):=&-\sum_{i=1}^n\dfrac{\partial}{\partial x_i}\left(\left|\dfrac{\partial u}{\partial x_i}\right|^{p-2}\dfrac{\partial u}{\partial x_i}\right)\\&+\bm{b}_1(\bm{x})\cdot\nabla u(\bm{x})+\alpha_0(\bm{x})|u(\bm{x})|^{p-2}u(\bm{x}),\end{aligned}
\end{equation}
其中 $p\in(1,\infty)$, 且 $\bm{b}_1$ 和 $\alpha_0$ 有界且连续. 注意, 除低阶项的区别外, 这与 \eqref{eq:anisotropic-p-power-divergence-example} 定义的算子相同. 请读者证明, 存在常数 $\widetilde C>0$, 使得若
\begin{equation}\label{eq:cv-quasilinear-negative-zero-order-coefficient-bound}
\alpha_0(\bm{x})>-\widetilde C,\quad\bm{x}\in\Omega,
\end{equation}
则假设 \ref{hyp:cv-pde-spatial-continuity}--\ref{hyp:cv-pde-uniform-gradient-coercivity} 成立 (Problem~\ref{ex:cv-pde-drift-example-hypotheses}). 由 Theorem~\ref{thm:cv-quasilinear-dirichlet-existence}, 有下面的存在性结果.
\end{Example}

\begin{Theorem}
\label{thm:cv-quasilinear-anisotropic-drift-existence}
设非线性二阶偏微分算子 $\widetilde A$ 由 \eqref{eq:cv-quasilinear-anisotropic-drift-example} 定义. 则对每个 $f\in W^{-1,q}(\Omega)$, 方程
\begin{equation}\label{eq:cv-quasilinear-anisotropic-drift-existence-equation}
\widetilde A(u)=f.
\end{equation}
存在弱解 $u\in W_0^{1,p}$.
\end{Theorem}

\Needspace{5\baselineskip}\textbf{Problems}

\begin{enumerate}[label=\thesection.\arabic*.,ref=\thesection.\arabic*,start=16,series=book-problems]

\item\label{ex:hemicontinuity-not-continuity-origin} 若对每个 $v,w\in X$,
\begin{equation}\label{eq:hemicontinuity-scalar-pairing-along-line}
\mathbb{R}\ni t\mapsto(T(u+tv),w)\in\mathbb{R}
\end{equation}
都连续, 则称映射 $T:X\to X^*$ 在 $u\in X$ 处是\textbf{半连续的}.\par 求一个在原点处半连续但不连续的函数 $f:\mathbb{R}^2\to\mathbb{R}^2$.

\item\label{ex:pseudo-monotone-surjectivity-proof} 证明 Theorem~\ref{thm:pseudo-monotone-surjectivity}.

\item\label{ex:browder-minty-hemicontinuity-surjectivity} 证明, 若将连续性假设替换为半连续性, 则 Theorem~\ref{thm:browder-minty-surjectivity} 仍成立.

\item\label{ex:pseudo-monotone-hemicontinuity-surjectivity} 证明, 若将连续性假设替换为半连续性, 则 Theorem~\ref{thm:pseudo-monotone-surjectivity} 仍成立.

\item\label{ex:bounded-monotone-implies-pseudo-monotone} 证明有界的单调算子为伪单调算子.

\item\label{ex:pseudo-monotone-strong-to-weak-continuity} 证明, 若 $T$ 为伪单调算子, 且 $u_j\to\overline u$ 在 $X$ 中 (强) 收敛, 则 $T(u_j)\rightharpoonup T(\overline u)$ 在 $X^*$ 中成立.

\item\label{ex:cv-type-operator-hemicontinuity} 证明变分法型算子为半连续的. (因此, 由 Problem~\ref{ex:pseudo-monotone-hemicontinuity-surjectivity}, 可去掉 Corollary~\ref{cor:cv-type-surjectivity} 中的连续性假设, 而结论仍成立.)

\item\label{ex:cv-pde-vanishing-weak-test-pairing} 设 $u_j\rightharpoonup\overline u$ 在 $W_0^{1,p}(\Omega)$ 中成立, 且 $\widehat T(u_j,v)\rightharpoonup\psi$ 在 $W^{-1,q}(\Omega)$ 中成立. 验证 \eqref{eq:cv-pde-vanishing-weak-test-pairing-exercise}.

\item\label{ex:cv-pde-induced-operator-coercivity-proof} 证明 Lemma~\ref{lem:cv-pde-induced-operator-coercivity}.

\item\label{ex:cv-pde-drift-example-hypotheses} 证明存在 $\widetilde C>0$, 使得若 \eqref{eq:cv-quasilinear-negative-zero-order-coefficient-bound} 成立, 则 \eqref{eq:cv-quasilinear-anisotropic-drift-example} 定义的拟线性微分算子 $\widetilde A$ 满足假设 \ref{hyp:cv-pde-spatial-continuity}--\ref{hyp:cv-pde-uniform-gradient-coercivity}. 指出假设 \ref{hyp:monotone-pde-measurability}--\ref{hyp:monotone-pde-growth} 中哪些对该算子\emph{不}成立.

\end{enumerate}
